\section{反標準orbitの完全部分圏と残余特異性}\label{sec:orbit-core}
本節は第\ref{sec:coherent-contraction}節の標準弱楕円gluingを固定する．
$S$は滑らかな有理弱del Pezzo曲面，$E\in|-K_S|$は滑らかな楕円曲線，
$q\in\Pic^0(E)$は任意とする．
$\Phi_q:\cD_q\to\Db(\qgr R_q)$を定理\ref{thm:weak-coherent-sections}の
自然関手，$\mathcal N_q=\ker\Phi_q$，
$\mathcal Q_q=(\cD_q/\mathcal N_q)^\natural$とする．
反標準orbitから生成される充満なthick部分圏を
\[
 \mathcal T_q=\thick_{\cD_q}\{s_q^ne:n\in\mathbb Z\},\qquad
 \mathcal P_q=\thick_{\Db(\qgr R_q)}\{\pi R_q(n):n\in\mathbb Z\}
\]
と書く．後者を\emph{orbit-perfect core}と呼ぶ．
これは非可換$qgr$のすべてのperfect対象を内在的に特徴付けた
という意味の名称ではない．

\begin{lemma}[弱gluingでのorbit間Ext]\label{lem:weak-orbit-ext}
全整数$a,b$について，sourceとtargetの両方で，
$s_q^ae,s_q^be$と$\pi R_q(a),\pi R_q(b)$の間のExtはそれぞれ
\[
\Ext^j(-,-)=
\begin{cases}
(R_q)_{b-a}&j=0,\ b\ge a,\\
(R_q)_{a-b-1}^{*}&j=2,\ b<a,\\
0&\text{それ以外}
\end{cases}
\]
となる．さらに，これらの同定の下で自然関手$\Phi_q$が誘導する写像は
全次数で同型である．
\end{lemma}
\begin{proof}
source側の計算は補題\ref{lem:orbit}の証明と同じ境界triangleから従う．
そこで使う入力は$e$のexceptionality，$s_q^{-1}\to\mathrm{id}$の
自然変換，境界での$\Psi_q^n\OO_E$の次数が$nd>0$であること，
および滑らかなgluingのSerre双対性であり，$-K_S$のamplenessは不要である．
正twistの$H^1$消滅により正の相対次数の高次Extが消え，
負の相対次数はSerre双対性で正側に帰着する．

target側は$R_q$が楕円座標環$B(E,\Psi_q)$の中心正則Rees変形であることと
Rains\cite[Corollary 3.8]{Filtered}から得られる．
この既知の計算はgenuineなblow-upの存在を仮定しない．
次数$0$の写像は$\Phi_q(s_q^ae)=\pi R_q(a)$と積の自然性から同型である．
次数$2$の同定には境界$\Ext_E^1(\OO_E,\OO_E)=H^1(E,\OO_E)$へのtraceを使う．
sourceとtargetで境界triangleと中心元$g$が同じなので，
$\Phi_q$はこのtraceを保存する．
特に$\Ext^2(e,s_q^{-1}e)\simeq H^1(E,\OO_E)$での写像は恒等で，
Yoneda積と非退化なSerre pairingから全ての負twistについて同型となる．
この比較は定理\ref{thm:derived}の完全忠実性証明のうち
\emph{orbit間の計算}だけを再利用し，Ind圏でのorbit生成は使っていない．
\end{proof}

\begin{theorem}[orbit-perfect core上の同値]\label{thm:orbit-core-equivalence}
任意の$q$について，$\Phi_q$は
\[
\boxed{\Phi_q|_{\mathcal T_q}:\mathcal T_q
  \xrightarrow{\sim}\mathcal P_q}
\]
というdg同値を与える．
さらに$\mathcal T_q$は$\mathcal N_q$と両側直交し，
商関手$\ell_q:\cD_q\to\mathcal Q_q$は$\mathcal T_q$上で完全忠実である．
従って$\mathcal P_q$は$\mathcal Q_q$内の
$\overline{\mathcal T}_q:=\ell_q(\mathcal T_q)$とdg同値である．
\end{theorem}
\begin{proof}
前補題はorbitの生成対象の間で$\Phi_q$が完全忠実であることを与える．
両変数について有限直和，shift，cone，直和因子を取れば
thick closure上でも完全忠実となり，像は定義より$\mathcal P_q$である．
$N\in\mathcal N_q$なら$\RHom(s_q^ne,N)=0$が全$n$で成立する．
定理\ref{thm:weak-orbit-kernel}より$s_q^{-1}N\simeq N$なので，
Serre双対性により
\[
 \RHom(N,s_q^ne)
 \simeq\RHom(s_q^ne,s_q^{-1}N[2])^{*}=0.
\]
したがって$\mathcal T_q$と$\mathcal N_q$は両側直交する．
Drinfeld商では核を経由するHomの補正項が消えるので，
$\ell_q$は$\mathcal T_q$上で完全忠実である．
\end{proof}

\begin{theorem}[零指数壁における残余特異性]\label{thm:residual-singularity}
\[
 \mathfrak S_q=
 \bigl(\mathcal Q_q/\overline{\mathcal T}_q\bigr)^\natural
\]
をorbit-perfect coreを除いた残余dg商とする．
\begin{enumerate}
\item $S$がgenuineでないなら，任意の$q$について
$\overline{\mathcal T}_q\subsetneq\mathcal Q_q$であり，
$\mathfrak S_q\ne0$である．
\item $q=\OO_E$では自然なdg同値
\[
 \boxed{\overline{\mathcal T}_{\OO_E}\simeq\Perf(\overline S),
 \qquad
 \mathfrak S_{\OO_E}\simeq D_{\rm sg}(\overline S)^\natural}
\]
が成立する．
\item 自然関手$\overline\Phi_q:\mathcal Q_q\to\Db(\qgr R_q)$は
$\overline{\mathcal T}_q$を$\mathcal P_q$に写すので，
\[
 \mathfrak S_q\longrightarrow
 \bigl(\Db(\qgr R_q)/\mathcal P_q\bigr)^\natural
\]
という関手を誘導する．$q=\OO_E$ではこれは上の古典的な
特異性圏の同値である．
\end{enumerate}
\end{theorem}
\begin{proof}
(1)は高次Extの障害による．$\cD_q$はproperなので，
$\mathcal T_q$の任意の二対象間の全次数のHomの総次元は有限である．
前定理により$\overline{\mathcal T}_q$も同じ性質を持つ．
一方，系\ref{cor:singular-gluing-unbounded}または
定理\ref{thm:duval-high-ext}から，非genuineな$\mathcal Q_q$には
無限個の次数で自己Extが非零となる対象がある．
従って$\overline{\mathcal T}_q$は$\mathcal Q_q$全体でない．
これはthick部分圏なのでVerdier商は非零である．

(2)では$\mathcal Q_{\OO_E}\simeq\Db(\coh\overline S)$が
定理\ref{thm:singular-gluing-target}で成立する．
反標準orbitの像は$\OO_{\overline S}(nH)$に等しい．
$H=-K_{\overline S}$はampleなので，その全整数twistは
$D_{\mathrm{qc}}(\overline S)$を生成する
（Stacks Project\cite{Stacks}，Tag 0BQQ）．
それらはcompactで，compact対象は$\Perf(\overline S)$だから，
そのthick closureは$\Perf(\overline S)$である．
(3)は関手の普遍性による．古典的$q=\OO_E$では
$\Phi_q$は$Rf_*$に一致し，対応する商も恒等となる．
\end{proof}
\begin{remark}[何を解決していないか]
$\mathfrak S_q$を構成したことは，
非自明$q$における
$\overline\Phi_q:\overline{\cD}_q\to\Db(\qgr R_q)$
の同値を証明するものではない．
特に$\mathcal Q_q$と$\Db(\qgr R_q)$が
同じorbit-perfect coreを持つとしても，
その補充部分やgluing bimoduleが一致するとは限らない．
零指数壁での障害は，単にorbit上のExtを比較するだけでは検出できず，
特異点に対応する非perfect対象とその高次Yoneda積を
標的で具体的に追う必要がある．
\end{remark}
